% =============================================================
% VI. A FINNISH MIRACLE THAT NEVER WAS
% =============================================================
\section{A Finnish miracle that never was}
\label{sec:finland}

The comparison with Finland has shaped the way Estonians and their historians think about the twentieth century. The two peoples speak related languages, were shaped by Lutheran schooling and Swedish law, emancipated their peasantries early and built co-operative agrarian economies after 1918; in the Soviet period Finland, visible on television from Tallinn, became `proof of what Estonia would have been if the country had not ended up under the Soviet occupation'.\footnote{\autocite[83]{lepppantti2013}.} The scholarly version of the argument has two steps. The first is that the two countries stood at the same level before 1940: Dellenbrant wrote that their levels of development were `close to equal around 1940', Soviet-era comparisons of real wages put Estonian workers ahead of Finnish ones as late as 1966, and the White Book's estimate of the losses of occupation rests on the assumption that Estonia would have followed the Finnish path.\footnote{\autocite[206]{dellenbrant1992}; \autocite[251, 257]{klesment2009}; \autocite[141--74]{kukk2005}.} The second is that the gap observed in 1991, when Estonian living standards were put at 40 per cent of Finland's, measures the cost of Soviet rule.\footnote{\autocite[206]{dellenbrant1992}.} Sceptics have questioned both steps. Kõll argued in 1994 that the belief in Baltic--Nordic parity before the war was `misleading'; an Estonian economist of the 1930s had already written that `Denmark, Sweden, Finland, and even Latvia are ahead of us regarding their wealth'; and Norkus, working from the benchmark estimates, concluded that the gap `cannot only be considered as the legacy of the Soviet occupation'.\footnote{\autocite[201]{koll1994gap}; \autocite[131]{klesment2005}; \autocite[18]{norkus2016}; \autocite[154]{norkus_restorations}.} The new series allows the two steps to be examined year by year.

Table~\ref{tab:relative} gives Estonian income per head as a percentage of Finland's at the benchmark years, and Table~\ref{tab:growth} the growth rates between them. In 1913 the two countries were indeed close: Estonia stood at \pctfin{1913} per cent of the Finnish level. In 1919, after wars that had damaged both, it stood at \pctfin{1919} per cent. The parity did not survive the 1920s. Finland grew by \grx{Finland}{1919}{1929} per cent a year to 1929 and Estonia by \gr{1919}{1929} per cent, and by 1929 Estonia was at \pctfin{1929} per cent of the Finnish level; the 1930s, in which Estonia's growth of \gr{1929}{1938} per cent a year was close to Finland's \grx{Finland}{1929}{1938} per cent, left the ratio at \pctfin{1938} per cent in 1938. Clark's 90 per cent, which Valge adopted and which has anchored the literature since, is thus too high by a sixth; Bocharnikova's finding that the purchasing power of wages in Estonia was only 4 per cent below Finland's in 1938 refers to urban wages alone and says nothing about the nearly two-thirds of the Estonian labour force on the land.\footnote{\autocite[81--2]{bocharnikova2021}.} The difference in 1938 was not one of social capability. Table~\ref{tab:structure} shows Estonian literacy above Finnish, Estonian university enrolment per head well above Finnish, and the same share of the labour force in agriculture; Table~\ref{tab:trade} shows both economies open, though Finland more so. The difference is better explained by resources and markets. Finland's forests, sawmills and pulp mills gave it an export staple for which Western demand recovered quickly after 1920 and which carried a growing engineering sector with it; Estonia's butter and bacon faced Danish competition and British quotas, and its industry had lost the only large market it had ever served.\footnote{\autocite[184, 198--200]{finlandeh01}; \autocite[193--210]{koll1994gap}.}

The second step of the argument fares no better, though for a different reason. Soviet rule did open a gap, but the timing is not what the argument supposes. Finland lost the Second World War and paid reparations; Estonia was fought over twice and collectivised. In 1950 the ratio stood at \pctfin{1950} per cent, higher than in 1938, and in 1968, after the two decades of Soviet industrialisation examined above, at \pctfin{1968} per cent, the highest level since 1919. Measured against Finland, the first half of Soviet rule was a period of convergence. The divergence that the literature attributes to occupation is a phenomenon of its second half: between 1968 and 1989 Finland grew by \grx{Finland}{1968}{1989} per cent a year, half as fast again as the Western core, and Estonia by \gr{1968}{1989} per cent, so that the ratio fell to \pctfin{1973} per cent in 1973 and \pctfin{1989} per cent in 1989. Those who lived through it dated the divergence correctly: Dellenbrant, drawing on Kukk's calculations, placed its beginning in the late 1960s.\footnote{\autocite[205]{dellenbrant1992}.} What the new series adds is the magnitude. Expressed in logarithms, the Estonian shortfall against Finland was 0.29 in 1938 and 0.63 in 1989; almost half of the gap at the end of Soviet rule had opened before it began, and the whole of the increase after 1938 occurred after 1968.

Two conclusions follow. The first concerns the use of Finland as a yardstick. Finland is a revealing comparator for Estonia, because it isolates the effect of institutions on two economies with similar endowments of human capital. But it is not a measure of what Estonia lost. The two economies were already a quarter apart in 1938, and Finland's growth after 1968, \grx{Finland}{1968}{1989} per cent a year, was exceptional by the standards of the Western core (\grx{Western}{1968}{1989} per cent) and of the Nordic countries (\grx{Nordics}{1968}{1989} per cent), driven by a forest-products and engineering complex that Estonia never possessed and by access to a Soviet market on favourable bilateral terms that Estonia, inside the Union, could not exploit.\footnote{\autocite[184--200]{finlandeh01}; \autocite[23--4]{norkus2016}.} The Western periphery as a whole, the group of poorer Western economies to which Estonia belonged by income in 1938, grew by \grx{Periphery}{1968}{1989} per cent a year over the same years (Table~\ref{tab:growth}); its Iberian and Greek members, starting from levels below Estonia's of 1950, grew as fast as Finland, while Ireland grew more slowly. Measured against the core or the periphery rather than against Finland, Estonia's shortfall in 1968--89 is between one and a half and two points a year rather than more than two and a half: still large, but a different magnitude from the one the White Book implies, and one that cannot be read off the Finnish series alone. The second concerns the Soviet record itself. The years in which Estonia kept pace with Finland were the years of forced industrialisation, high investment and immigration; the years in which it fell behind were those in which the planned economy, having exhausted the sources of extensive growth, could neither raise productivity nor change its structure. The gap with Finland in 1991 is therefore likely to overstate the cost of occupation. The part of it that can be charged to the system is the stagnation of the last two decades, and the measure of that cost is Estonia's shortfall over those years against the Western economies of its own level rather than against Finland: a less comforting conclusion than the one the White Book drew, and a more exact one.
